\documentclass[12pt]{article}
\usepackage[utf8]{inputenc}
\usepackage[T1]{fontenc}
\usepackage{hyperref}
\title{Toma Language Resources for K-12 Curriculum Development}
\author{Compiled by Claude Sonnet 5 (claude-sonnet-5) for Andr\'as Kornai}
\date{}
\begin{document}
\maketitle

\section{Introduction}
Toma (ISO 639-3: tod; Glottocode: toma1245) is the name for the Guinean varieties of Loma, a Mande language. Omniglot puts speakers in Guinea at about 261,000. The language is written in a Latin-based alphabet. A syllabary from the 1930s fell out of use, and revival efforts are reported. This search found little classroom material.

\section{Primary Educational Resources}
Omniglot mentions some literature and a literacy programme, but I found no named primers or textbooks. See \url{https://omniglot.com/writing/loma.htm}.

\section{Linguistic and Reference Materials}
The Wikipedia article on Loma covers Toma in Guinea: \url{https://en.wikipedia.org/wiki/Loma_language}. A Unicode proposal on encoding the Loma script cites alphabet sources, including a 1993 Loma alphabet introduction by Justin Kaliva Koïvogui and an undated ``Alphabet Toma Macenta'' by Foromo Guilavogui: \url{https://www.unicode.org/L2/L2017/17059-n4786-loma.pdf}. I found no dictionary.

\section{Religious and Community Resources}
Global Recordings lists audio Bible content in Toma: \url{https://www.globalrecordings.net/en/language/tod}.

\section{Digital and Audio Resources}
Endangered Alphabets describes the syllabary revival: \url{https://www.endangeredalphabets.net/loma}. See also \url{https://digitalorientalist.com/2021/03/19/advancing-on-the-encoding-of-loma/}.

\section{Curriculum Development Context}
A Latin orthography and a literacy programme exist. Teaching materials, dictionaries and grammars were not found in this search and need a follow-up.
\end{document}